% Geometry, mesh and result products in the current workflow.
\section{Geometry, mesh and results}

\begin{frame}[fragile]{Build the shape, then the mesh}
\begin{lstlisting}
cfd-geom validate model.toml
cfd-geom run model.toml
cfd-mesh airfoil-2d --naca 0012 --reynolds 6000000
\end{lstlisting}
\begin{itemize}\small
  \item \code{cfd-geom} builds parametric geometry and records its stages in
        \code{inputs/geometry/<name>/}, with \code{geometry.json}.
  \item \code{cfd-mesh} generates and transforms meshes in \code{inputs/mesh/},
        with provenance sidecars. The airfoil command writes SU2 and gmsh files.
  \item A 2D SU2 or OpenFOAM row that sweeps \code{ALPHA} can become a section
        polar: \code{polars/P<POL>.txt}, readable by XFOIL and QPROP.
\end{itemize}
\end{frame}

\begin{frame}{From a saved solution to a recorded result}
\begin{itemize}\small
  \item \code{RESTART} in a row names a recorded source solution;
        \code{run --restart} continues each point from its own solution.
        SU2, OpenFOAM and PyFR support both. Earlier evidence is archived first.
  \item The host guard is \textbf{off by default}. Add
        \code{--one-run-at-a-time}, or set \code{one_run_at_a_time = true}
        under \code{[solver]} in \code{init.toml}, to reserve the host.
  \item In \code{VAR_NAMES_VALUES}, \code{hpc: alpha / cluster: west}
        selects a profile and cluster per row. \code{hpc:} overrides
        \code{--hpc}; \code{--local} ignores both row selections.
  \item \code{collect} writes \code{campaign_summary.json}, family tables such as
        \code{FAM-ACOEF-P<POL>.csv}, group tables and \code{products.json}.
        Setup \code{[groups]} boundary forces stay in the summary, separately
        from the family-based group CSV files. Guide 2 gives the file map.
\end{itemize}
\end{frame}
